\section{Fundamentação teórica}

\subsection{Inferência de ensembles de árvores em ONNX Runtime}

XGBoost \cite{chen2016} é um algoritmo de \textit{gradient boosting} para
\textit{ensembles} de árvores de decisão que domina problemas de classificação
e regressão sobre dados tabulares densos, frequentemente superando redes
neurais em conjuntos representativos \cite{grinsztajn2022}. O modelo treinado
neste trabalho contém 500 árvores binárias de profundidade máxima 8, ou seja,
aproximadamente 128 mil nós e até 4 mil comparações por inferência sobre uma
única linha.

Para servir o modelo em produção, ele é exportado para o formato ONNX
\cite{onnxruntime2024}, que define o operador
\textit{TreeEnsembleClassifier} sob o domínio \texttt{ai.onnx.ml}. O operador
recebe um tensor de \textit{features} de forma $[N, 28]$ em
\texttt{float32}, percorre cada árvore de forma independente e acumula as
contribuições das folhas, produzindo um tensor de probabilidades de forma
$[N, 2]$ \cite{onnx_treeensemble}. A implementação no ONNX Runtime usa um
arranjo plano de nós, paralelizando entre árvores via \textit{thread pool}
sempre que \texttt{intra\_op\_num\_threads} é maior que um, e libera o
\textit{Global Interpreter Lock} (GIL) durante a chamada nativa de
\texttt{session.run}.

O custo intrínseco da inferência depende portanto da profundidade média, do
número de árvores, do tamanho do conjunto de trabalho na cache da CPU e da
qualidade do preditor de saltos. Asadi et al. \cite{asadi2014} estabelecem
um modelo no qual cada comparação custa de 2 a 5 nanossegundos em
processadores modernos com arranjos amigáveis à predição de salto. Para o
modelo aqui empregado, isso prediz um piso teórico entre 8 e 20
microssegundos por inferência sobre uma linha. Sieboehm
\cite{sieboehm2021} mediu o ONNX Runtime em um modelo \textit{LightGBM}
comparável, em \textit{hardware} Intel Haswell, obtendo 11,00 microssegundos
por linha, o que corrobora a ordem de grandeza esperada. A consequência
prática é que, para esse regime, o sobrecusto da camada de \textit{serving}
ao redor da inferência (parsing, despacho HTTP, marshaling de tensores)
domina a latência ponta a ponta.

\subsection{Latência de cauda e Coordinated Omission}

A literatura de sistemas estabelece o percentil 99 (P99) da distribuição de
latência como a métrica operacional dominante para serviços interativos,
porque a mediana subestima as anomalias raras que afetam a experiência do
usuário \cite{dean2013}. Em uma distribuição típica de \textit{serving} sob
carga concorrente, a cauda é dominada por pausas do \textit{garbage
collector}, contenção de recursos, \textit{cache misses} e variação de
escalonamento, eventos que não desaparecem por aumento do número de
amostras e que tendem a se amplificar quando o sistema se aproxima da
saturação.

Medir a cauda corretamente exige um gerador de carga de modelo aberto. Em um
gerador de modelo fechado, cada nova requisição espera a resposta da
anterior, de modo que quando o servidor engasga, a taxa de envio cai e a
cauda observada subestima a cauda real, fenômeno descrito por Tene como
\textit{Coordinated Omission} \cite{tene2015}. Em um gerador de modelo
aberto, as requisições são despachadas a intervalos pré-definidos
$1/\text{RPS}$, independentemente do número de requisições em voo. Quando o
sistema se torna incapaz de atender à taxa de chegada, a discrepância
aparece como crescimento de fila e latência, e não como redução da carga
injetada. Vegeta \cite{vegeta2024} é um gerador de modelo aberto que registra
a latência percebida pelo cliente por requisição em formato binário, do qual
percentis arbitrários, função de distribuição complementar acumulada (CCDF)
e vazão útil sob SLA podem ser extraídos em pós-processamento.

\subsection{Adaptive batching e o modelo de custo do despacho}

\textit{Adaptive batching} é uma técnica em que requisições concorrentes que
chegam ao servidor são agrupadas em lotes antes da inferência, com o objetivo
de amortizar um custo fixo por chamada $k$ entre $B$ requisições. O ganho
líquido por requisição é aproximadamente $k(1 - 1/B) - \tau_{\text{wait}}$,
onde $\tau_{\text{wait}}$ é o tempo médio que cada requisição passa
aguardando o fechamento do lote. Para modelos com $k$ elevado, como redes
profundas em GPU (lançamento de \textit{kernel}, materialização de tensores,
\textit{boilerplate} de \textit{Python}), o lote de $B = 16$ pode reduzir o
custo por requisição em uma ordem de grandeza
\cite{nakandala2020}. Para inferência de árvores em CPU, onde $k$ é da ordem
de 14 microssegundos, o sobrecusto de fila e do mecanismo de despacho
sobrepuja qualquer amortização possível.

O \textit{dispatcher} do BentoML é implementado como uma única corrotina
\texttt{controller()} por \textit{worker} \cite{bentoml_dispatcher}. Essa
corrotina drena a fila de requisições, decide o momento do disparo de cada
lote a partir de estimadores de profundidade de fila e do tempo de
serviço, e roteia os resultados de volta para as \textit{futures}
correspondentes. O parâmetro \texttt{max\_latency\_ms} é interpretado como um
prazo rígido para a conclusão da requisição: quando a latência estimada
ultrapassa o limite, a requisição é cancelada e a resposta retorna como HTTP
503. Essa interpretação difere da intuição de janela de espera a ser
minimizada e tem consequências diretas para a configuração do servidor sob
carga.

\subsection{Métodos estatísticos para comparação entre pipelines}

Cada célula da varredura matricial é executada três vezes, e as métricas
reportadas no nível da célula são as medianas entre as repetições. Para
acompanhar a variabilidade entre repetições, dois resumos são computados.
O primeiro é o coeficiente de variação $\text{CoV} = \sigma/\mu$, calculado
sobre as três medidas de P50, P95 e P99, e usado como indicador de
estabilidade entre execuções. O segundo é o intervalo de confiança bootstrap
de 95\% sobre cada percentil reportado, obtido por reamostragem com
reposição de 1.000 réplicas das três medidas e tomando os quantis 2,5\% e
97,5\% da distribuição resultante \cite{efron1979}. Com $n = 3$ repetições os
intervalos são largos em termos absolutos, e por isso as comparações
entre pipelines são apoiadas pela magnitude da diferença, não pela estreitez
dos intervalos.

Como o BentoML descarta requisições por prazo expirado (HTTP 503) enquanto
o FastAPI as enfileira como latência adicional, comparar taxas de erro
diretamente seria injusto: um 503 em 5\,ms e uma resposta 200 em 2\,s sob
um SLA de 50\,ms são ambos falhas operacionais idênticas do ponto de vista
do consumidor. A métrica adotada para a comparação principal é a vazão útil
sob SLA, definida como a fração de requisições oferecidas que retornaram
status 2xx com latência menor ou igual a um limite (50 ou 100\,ms), divida
pelo total oferecido. Essa formulação trata os dois modos de falha de forma
simétrica e produz uma única figura que pode ser comparada entre pipelines
e células.

